\section{Ayokha --- ``The Monsoon Throne''}
\label{chap:ayokha}

\begin{quote}
  \textbf{God-King's Astrologer (Elite):}  
  ``The monsoon does not come because we pray. It comes because the stars align. I read the stars. I align them. Without me, the rains would fail, the rice would burn, and the kingdom would fall. I am not a servant of the God-King. I am his monsoon.''
\end{quote}

\begin{quote}
  \textbf{River Merchant (Commoner):}  
  ``My boat has more secrets than cargo. The lower deck is for silk. The upper deck is for lies. The tiller is for bribes. I have carried nine spies this year. The ninth was the God-King's own daughter.''
\end{quote}

\begin{tcolorbox}[colback=black!3,colframe=blue!60!black,title={Theme \& Atmosphere}]
Ayokha is the land of the Monsoon Throne, where the God-King rules by the rhythm of the rains. Rivers flood, rice terraces glow green, and the wharves are crowded with ships from Dhahara, Sihai, and distant isles. The Lethai-ar witness every contract; the spirit-mediums speak for the nats of the ford. Power here is measured in monsoon schedules, silk vouchers, and the favour of the God-King's astrologer. The ninth taboo is strong: never pour the ninth cup, never count the ninth wave, never speak the ninth name. The Hollow watches from the flooded fields.

\vspace{2mm}
\textbf{What you notice first:} The smell of wet earth and fermenting rice. The way every bridge has a spirit shrine at its foot. The constant sound of water — dripping, flowing, raining.

\textbf{The one rule that kills newcomers:} Never refuse a gift of betel leaf. It is an oath of hospitality. Refusing it is an insult that can end a bloodline.
\end{tcolorbox}

\subsection*{Regional Mood: The Weight of the Monsoon}

Power here is measured in monsoon predictions, silk bales, and the number of spirits placated. The God-King rules, but the astrologer sets the calendar, the Lethai-ar witness the contracts, and the spirit-mediums negotiate with the nats. A merchant who ignores a nat's demand will find his boat sunk. A farmer who disrespects the monsoon will find his fields barren.

\textbf{Starting Location:} A river market on stilts, where boats load silk and spices. A Lethai-ar Vigil holds a green lamp: ``The God-King's astrologer has predicted a late monsoon. The merchants are panicking. The rice will fail unless someone sacrifices to the river spirit. You are someone. Go.''

\subsection*{Prominent NPCs of Ayokha}
\label{sec:ayokha-npcs}

\begin{longtable}{|p{0.2\textwidth}|p{0.25\textwidth}|p{0.3\textwidth}|p{0.15\textwidth}|}
\hline
\textbf{Name} & \textbf{Role/Title} & \textbf{Motivation} & \textbf{Rumoured Quirk} \\
\hline
God-King's Astrologer & Royal astronomer & Predict the monsoon before the people starve & He has not seen the sun in thirty years. He works by candlelight. \\
\hline
Vigil Lady Suyin & Lethai-ar witness-master & Keep the river's ledger free of forgeries & She has cut nine witness cords. The ninth was her own. \\
\hline
River Merchant & Silk trader & Smuggle contraband under the God-King's tariffs & She has a hidden deck that only appears at low tide. \\
\hline
Spirit-Medium & Nat negotiator & Speak for the spirit of the ford & She carries a bowl of water; the spirit speaks through ripples. \\
\hline
The Drowned Nat (Ghost) & A spirit of a broken oath & Demand payment before the monsoon fails & It appears as a wet figure in torn robes. It never blinks. \\
\hline
\end{longtable}

\subsection*{Names of Ayokha}
\label{sec:ayokha-names}

Inspired by the monsoon kingdoms: soft, flowing consonants, honorifics (Lady, Master). Surnames denote a river, a wharf, or a spirit.

\begin{longtable}{|p{0.25\textwidth}|p{0.25\textwidth}|p{0.2\textwidth}|p{0.2\textwidth}|}
\hline
\textbf{First Names (Male)} & \textbf{First Names (Female)} & \textbf{Clan/Epithet} & \textbf{Monsoon-Name} \\
\hline
Hiro & Suyin & River-watch & \textit{the Uncut} \\
\hline
Tama & Mira & Silk-hand & \textit{Ninth Knot} \\
\hline
Kaito & Yuki & Lamp-bearer & \textit{Green Flame} \\
\hline
Riku & Hana & Oath-strand & \textit{the Witness} \\
\hline
Sora & Nami & Mud-dyer & \textit{the Unforgotten} \\
\hline
\end{longtable}

\paragraph*{(Market/Temple/Bridge)} A river market on stilts; a step-pyramid temple with monsoon carvings; a bridge that appears only once a century.

\section*{Spades — Places}
\label{sec:ayokha-places}
\begin{enumerate}
\setcounter{enumi}{1}
\item \textbf{River Market (Stilts)} --- A floating market where goods are passed from boat to boat.
  \hfill \textit{The prices are written in water. They change with the tide.}
\item \textbf{Step-Pyramid Temple} --- A stone pyramid rising from the riverbank. Carvings of the first monsoon.
  \hfill \textit{The carvings weep water during the rainy season. The water is holy.}
\item \textbf{Bridge of Three Echoes} --- A wooden bridge that appears only once a century.
  \hfill \textit{Those who cross it see three versions of themselves: past, present, and future.}
\item \textbf{Silk Village (Mulberry Groves)} --- A village where silk is spun and cords are woven.
  \hfill \textit{The mulberry trees were planted by the first God-King. Their leaves are sacred.}
\item \textbf{Astrologer's Observatory} --- A roofless tower with a bronze armillary sphere.
  \hfill \textit{The sphere has nine rings. The ninth ring is broken. The astrologer refuses to fix it.}
\item \textbf{River Court (Floating)} --- A barge that serves as a mobile courthouse.
  \hfill \textit{The barge is anchored at a different spot each day. The river is the judge.}
\item \textbf{Spirit-Medium's Hut} --- A small shelter on the riverbank, surrounded by offering bowls.
  \hfill \textit{The bowls are filled with rice, flowers, and small coins. The spirits take the offerings at night.}
\item \textbf{Well of the Drowned Nat} --- A stone well where a spirit was trapped.
  \hfill \textit{The well is covered by a stone slab. The slab has nine holes. The spirit breathes through them.}
\item \textbf{God-King's Palace (Floating)} --- A wooden palace built on pontoons, anchored in the river.
  \hfill \textit{The palace rises and falls with the flood. The God-King never touches dry land.}
\item[J] \textbf{The Ninth Terrace} --- A rice terrace that is flooded only once every nine years.
  \hfill \textit{The rice grown there is for the God-King alone. Eating it is treason.}
\item[Q] \textbf{Monsoon Bell Tower} --- A stone tower with a single brass bell. The bell is rung at the first rain.
  \hfill \textit{The bell's tone determines the harvest. A high tone means abundance. A low tone means famine.}
\item[K] \textbf{The First Bridge (Myth)} --- Where the first God-King crossed the river and claimed the throne.
  \hfill \textit{The bridge is invisible. Only those who have witnessed nine true oaths can see it.}
\item[A] \textbf{The Monsoon's Eye} --- A still pool in the astrologer's observatory. It reflects the sky.
  \hfill \textit{The pool shows the future of the monsoon. The astrologer is the only one who can read it.}
\end{enumerate}

\paragraph*{(Merchant/Medium/Vigil)} River merchant with a boat of secrets; spirit-medium with a bowl of water; Lethai-ar Vigil with a green lamp.

\section*{Hearts — People & Factions}
\label{sec:ayokha-people}
\begin{enumerate}
\setcounter{enumi}{1}
\item \textbf{River Merchant} --- A woman with a brass earring and a hidden deck. She never sails without a bribe.
  \hfill \textit{``The tariff is negotiable,'' she says. ``My silence is not.''}
\item \textbf{Spirit-Medium} --- A woman with a bowl of water. She watches the ripples.
  \hfill \textit{``The nat says the river is angry,'' she says. ``Someone broke a promise.''}
\item \textbf{Vigil Lady Suyin} --- Grey-robed, a green lamp at her hip. She never blinks when witnessing an oath.
  \hfill \textit{``The cord will tighten,'' she says. ``I do not need to watch. The river watches for me.''}
\item \textbf{God-King's Astrologer} --- A pale man in silk robes, surrounded by candles.
  \hfill \textit{``The monsoon will be late,'' he says. ``The rice will fail unless...'' He trails off.}
\item \textbf{Rice Farmer} --- A woman with mud-stained feet and a straw hat. She leaves offerings at the terrace edge.
  \hfill \textit{``The spirits are hungry,'' she says. ``Feed them, or they feed on you.''}
\item \textbf{Silk Weaver} --- A man with thread between his fingers. He weaves witness cords for the Lethai-ar.
  \hfill \textit{``The cord must be perfect,'' he says. ``A single flaw can break an oath.''}
\item \textbf{Monsoon Bell Ringer} --- A young woman who rings the bell at the first rain. She has never been late.
  \hfill \textit{``The rain waits for me,'' she says. ``I do not wait for the rain.''}
\item \textbf{Drowned Nat's Ghost} --- A wet figure in torn robes, standing by the well.
  \hfill \textit{It asks: ``Who sealed me here?'' The answer is a name. The name is the God-King's.}
\item \textbf{God-King (Offstage)} --- Never seen, rarely heard. His decrees are delivered by astrologer.
  \hfill \textit{Rumour says he is a spirit. Rumour says he is a child. Rumour says he is the river itself.}
\item[J] \textbf{The First God-King's Ghost} --- A crowned skeleton in the palace's lower deck.
  \hfill \textit{He asks: ``Is the monsoon still on time?'' If you say no, he weeps.}
\item[Q] \textbf{Monsoon Spirit (Kami)} --- A whirlwind with eyes, seen over the river during storms.
  \hfill \textit{It asks: ``Who rang the bell?'' The answer must be the ringer's name.}
\item[K] \textbf{The Ninth Terrace's Keeper} --- An old woman who tends the forbidden rice.
  \hfill \textit{She has never eaten the rice. She has never been tempted.}
\item[A] \textbf{Kuva's Shadow (Hollow)} --- A vast darkness that covers the river at midnight.
  \hfill \textit{It asks: ``Who poured the ninth cup?'' The answer is a name. The name is the speaker's.}
\end{enumerate}

\paragraph*{(Flood/Rumor/Offense)} The river floods — a bridge is swept away; a rival merchant spreads rumors of cursed silk; the nat of the well is offended — no one may draw water.

\section*{Clubs — Complications/Threats}
\label{sec:ayokha-complications}
\begin{enumerate}
\setcounter{enumi}{1}
\item \textbf{River Floods} --- The monsoon rises. The bridge is gone. The wharf is isolated.
  \hfill \textit{The flood is not natural. A river-spirit is angry. An oath was broken upstream.}
\item \textbf{Rumour of Cursed Silk} --- A merchant claims that a bale of silk is cursed. Trade stops.
  \hfill \textit{The curse is a lie. The merchant wants to drive down prices.}
\item \textbf{Well-Spirit Offended} --- A traveller threw trash into the sacred well. The spirit refuses water.
  \hfill \textit{The spirit demands an apology. The apology must be an oath. The oath must be witnessed.}
\item \textbf{Monsoon Late} --- The astrologer's prediction fails. The rains do not come.
  \hfill \textit{The astrologer is a fraud. Or the spirits are angry. Either way, the rice is dying.}
\item \textbf{Witness Cord Cut (False)} --- Someone cuts a witness cord that was tied in good faith.
  \hfill \textit{The cutter is the Vigil's own apprentice. She thought the oath was false. It was true.}
\item \textbf{Silk Village Poisoned} --- Someone poisons the mulberry trees. The silkworms die.
  \hfill \textit{The poisoner is a rival merchant. He wants to monopolise the silk trade.}
\item \textbf{River-Spirit Demands a Story} --- The nat appears at the ford. It asks for a broken oath.
  \hfill \textit{It will not leave until someone confesses. The confession must be true.}
\item \textbf{God-King's False Oath} --- The king swears a vow he does not intend to keep. The Vigil knows.
  \hfill \textit{If she exposes him, the king will exile the Lethai-ar. If she stays silent, the cord tightens.}
\item \textbf{Astrologer's Betrayal} --- The astrologer sells the monsoon schedule to a foreign power.
  \hfill \textit{The foreign power is Ashaan. They want to time an invasion with the floods.}
\item[J] \textbf{The Ninth Cup Poured (By Mistake)} --- A stranger pours the ninth cup at a river market. The Hollow comes.
  \hfill \textit{The Hollow asks: ``Who poured this cup?'' The stranger points at the party.}
\item[Q] \textbf{First Bridge Appears} --- The mythic bridge manifests. Crossing it grants a new name.
  \hfill \textit{The new name is a debt. The debt is to the first God-King. He will collect.}
\item[K] \textbf{Monsoon Bell Cracked} --- The bell is struck too hard. It cracks. The tone is wrong.
  \hfill \textit{The bell-ringer is a spy. She cracked it on purpose. The monsoon will fail.}
\item[A] \textbf{The Monsoon Fails} --- The rains do not come for nine weeks. The river dries. The Hollow walks.
  \hfill \textit{The Hollow asks: ``Who failed to ring the bell?'' The answer is a name.}
\end{enumerate}

\paragraph*{(Schedule/Voucher/Cord)} Monsoon schedule (predicts safe sailing days); silk voucher (trade goods at 10\% profit); Lethai-ar witness cord (enforce any oath).

\section*{Diamonds — Rewards/Leverage}
\label{sec:ayokha-rewards}
\begin{enumerate}
\setcounter{enumi}{1}
\item \textbf{Monsoon Schedule (Astrologer's Copy)} --- A scroll of safe sailing days for the next season.
  \hfill \textit{The schedule is written on dried palm leaf. The leaf will crumble after use.}
\item \textbf{Silk Voucher (Merchant's Seal)} --- A stamped document. Trade goods at 10\% profit in any Ayokha port.
  \hfill \textit{The seal is a fish. The fish glows when the voucher is genuine.}
\item \textbf{Lethai-ar Witness Cord} --- A silk cord dyed with river mud. Enforce any oath.
  \hfill \textit{The cord tightens when the oath is broken. It will not stop until blood is drawn.}
\item \textbf{God-King's Blessing (Scroll)} --- A royal decree. One national favour in Ayokha.
  \hfill \textit{The scroll is sealed with river clay. The clay is from the palace's foundation.}
\item \textbf{River Merchant's Hidden Deck Key} --- A brass key. Access any smuggler's hold on the river.
  \hfill \textit{The key is warm. It glows when a customs boat is near.}
\item \textbf{Spirit-Medium's Bowl} --- A clay bowl. Fill it with water to speak with any river-spirit once.
  \hfill \textit{The spirit will answer one question. The answer will be a ripple.}
\item \textbf{Vigil's Lamp (Green Flame)} --- A small oil lamp. Gain +1 die to social rolls in any Lethai-ar court.
  \hfill \textit{The flame flickers when someone lies. It also flickers when you lie.}
\item \textbf{Silk Weaver's Needle} --- A needle carved from a mulberry branch. Tie a witness cord without cutting the silk.
  \hfill \textit{The needle can also cut a cord without making a sound. Use once.}
\item \textbf{Monsoon Bell Tone (Recording)} --- A clay cylinder that captured the bell's ring. Play it to summon a monsoon spirit.
  \hfill \textit{The spirit will come. It will demand an offering. The offering must be a promise.}
\item[J] \textbf{First Bridge Plank} --- A piece of the mythic bridge. Cross any river as if on solid ground once.
  \hfill \textit{The plank crumbles after use. The first God-King will remember.}
\item[Q] \textbf{Ninth Terrace Rice (Grain)} --- A single grain from the forbidden terrace. Eat it to gain a vision of the future.
  \hfill \textit{The vision is always a flood. The flood is always coming.}
\item[K] \textbf{Astrologer's Sphere Ring (Broken)} --- The ninth ring from the armillary sphere. Touch it to predict the monsoon's arrival.
  \hfill \textit{The ring is cold. It warms when the rains are nine days away.}
\item[A] \textbf{The Monsoon's Eye (Water)} --- A vial of water from the astrologer's pool. Drink it to see the future of the harvest.
  \hfill \textit{The water tastes of salt. The vision is always a famine. The famine is not inevitable.}
\end{enumerate}

\section*{Quick use notes}
\label{sec:ayokha-quick-use}
\begin{itemize}
\item Draw until you have all four suits. Highest rank sets main Clock.
\item Diamonds are codified outcomes (schedules, vouchers, cords) that change position rather than call for a roll.
\item If any A appears, echo \textbf{monsoon-omens}—water that flows upstream, a lamp that burns green without oil, a cord that ties itself.
\end{itemize}

\subsection*{Additional Features}
\label{sec:ayokha-features}

\begin{itemize}
\item \textbf{Monsoon Clock [6]:} Tracks the approach of the rainy season. When it fills, the monsoon arrives — all river travel becomes Desperate for 2 legs. The clock ticks each day during the dry season. The astrologer can predict the exact day.
\item \textbf{River Court Justice:} Disputes on the river are settled by a floating court. The judge is a Lethai-ar Vigil. The verdict is written on a clay tablet and thrown into the water. If the tablet floats, the verdict is true. If it sinks, the verdict is false — and the judge is drowned.
\item \textbf{Betel Leaf Hospitality:} Offering a betel leaf is an oath of hospitality. Refusing is a grave insult. Accepting binds both parties to non-violence for the duration of the visit.
\end{itemize}

\subsection*{Lore Echoes (Cross-Expansion Hooks)}
\begin{itemize}
  \item \textbf{Monsoon Trade and Dhahara:} Ayokha's monsoon winds carry ships to the Brass Coast. A delay in Ayokha means famine in Dhaharan ports.
  \item \textbf{Witness Cords and the Sidhi (Way of Silk):} Sidhi freedmen carry Lethai-ar witness cords as proof of emancipation. A cord from Vigil Suyin is worth more than a freedman's seal.
  \item \textbf{River Spirits and the Ukwe (Sekogo):} The Ayokhan river is said to flow from the Ukwe's roots. A spirit-medium who visits the Ukwe can read the river's memory without a context key.
\end{itemize}

\subsection*{Ayokha Superstitions (burn for +1 die once)}
\begin{itemize}
  \item \textbf{Bread to the River:} Crumble a crust into the current. Ignore the first \emph{River Floods} penalty.
  \item \textbf{Knot the Cord:} Tie a single knot in a witness cord before using it. Cancel \emph{Witness Cord Cut} or take –1 die on your next oath.
  \item \textbf{Name to the Well:} Whisper a dead Vigil's name into the well. Gain +1 Effect when witnessing an oath, but mark \emph{Obligation} to the Monsoon Spirit.
\end{itemize}

\begin{tcolorbox}[colback=black!3,colframe=black!40!white,title={Patronage \& Power}]
In Ayokha, power flows through the monsoon, the river, and the weight of a witnessed oath. The God-King rules, but the astrologer sets the calendar, the Lethai-ar hold the ledgers, and the spirit-mediums speak for the nats.

\textbf{For the GM:}  
Use the \textbf{Monsoon Clock} as a visible campaign timer. Let \textbf{River Court Justice} be a high-stakes social encounter. Enforce \textbf{Betel Leaf Hospitality} as a binding ritual — breaking it invites the Hollow.
\end{tcolorbox}
